\documentclass[11pt,a4paper]{article}
\usepackage{fontspec}
\usepackage{polyglossia}
\setmainlanguage{russian}
\setotherlanguage{english}
\setmainfont{Liberation Serif}
\usepackage[left=1.0cm,right=1.0cm,top=0.8cm,bottom=0.7cm]{geometry}
\usepackage{amsmath,amssymb}
\usepackage{array}
\usepackage{multirow}
\setlength{\parindent}{0pt}
\pagestyle{empty}

\newcounter{zad}

% ФОРМАТ. Макросы взяты из listok_22-09_podschety.tex без изменений.
% Поле ответа слева, условие справа; верх рамки поля стоит на одной
% горизонтальной линии с верхом строки условия (подгонка — \podnyat).
% Полоска под плюс ОДНА на задачу: плюс ставится за задачу целиком.
%
% ОТЛИЧИЯ ЭТОГО ЛИСТКА ОТ СПЕЦМАТОВСКИХ:
%   · принимается ПО РЕШЕНИЮ, но поле ответа есть — оно для САМОПРОВЕРКИ
%     школьника перед тем, как поднять руку;
%   · один вариант (варианты будут у контрольной, не здесь);
%   · значок только у звёздочек; уровни ○ и △ не используются;
%   · листок аскетичный: условия, поля, подпись. Всё остальное — устно.
%
% ПОРЯДОК ЗАДАЧ — по возрастанию сложности.
% Буква n стоит ТОЛЬКО там, где от неё зависит ТИП ответа: задача 4
% (кратность трём) и задача 6 (чётность). Везде остальное — конкретные
% числа, потому что там n меняет только величину ответа, а не его природу.
% В задаче 7 числа нет вовсе («несколько жителей») — и это тот же принцип:
% чётность там не дана, а ДОКАЗЫВАЕТСЯ, поэтому называть её нельзя.
%
% ЛИНИИ (каждая: опора внизу, вершина выше):
%   А — утверждения о самих себе:      задача 1 → задача 6
%   Б — конфигурация с параметром:     задача 3 → задача 4
%   В — доказать, что иначе нельзя:    задача 2 → задача 5 → задача 7

\newlength{\otvsh}\setlength{\otvsh}{3.5cm}
\newlength{\plsh}\setlength{\plsh}{0.7cm}
\newlength{\polesh}\setlength{\polesh}{4.6cm}
\newlength{\uslsh}\setlength{\uslsh}{13.6cm}
\newlength{\podnyat}\setlength{\podnyat}{1.1ex}
\newlength{\minpole}\setlength{\minpole}{1.3cm}

\newsavebox{\uslbox}
\newlength{\uslvys}
\newlength{\blh}

\newcommand{\calcblh}[1]{%
  \setlength{\blh}{\dimexpr(\uslvys-\arrayrulewidth)/#1-\arrayrulewidth\relax}}
\newcommand{\mt}[1]{\parbox[t][\blh][t]{\otvsh}{\small\bfseries #1\strut}}

\newcommand{\poleA}{\calcblh{1}\setlength{\tabcolsep}{3pt}%
  \begin{tabular}{|p{\otvsh}|p{\plsh}|}\hline \mt{} & \\ \hline\end{tabular}}
\newcommand{\poleB}[2]{\calcblh{2}\setlength{\tabcolsep}{3pt}%
  \begin{tabular}{|p{\otvsh}|p{\plsh}|}\hline
  \mt{#1} & \multirow{2}{*}{} \\ \cline{1-1} \mt{#2} & \\ \hline\end{tabular}}
\newcommand{\poleC}[3]{\calcblh{3}\setlength{\tabcolsep}{3pt}%
  \begin{tabular}{|p{\otvsh}|p{\plsh}|}\hline
  \mt{#1} & \multirow{3}{*}{} \\ \cline{1-1} \mt{#2} & \\ \cline{1-1}
  \mt{#3} & \\ \hline\end{tabular}}
\newcommand{\poleD}[4]{\calcblh{4}\setlength{\tabcolsep}{3pt}%
  \begin{tabular}{|p{\otvsh}|p{\plsh}|}\hline
  \mt{#1} & \multirow{4}{*}{} \\ \cline{1-1} \mt{#2} & \\ \cline{1-1}
  \mt{#3} & \\ \cline{1-1} \mt{#4} & \\ \hline\end{tabular}}

\newlength{\mezhzad}\setlength{\mezhzad}{14pt}
\newcommand{\znak}[1]{\ifx&#1&\else$#1$\ \fi}

\newcommand{\zad}[3]{%
  \stepcounter{zad}%
  \sbox{\uslbox}{\parbox[t]{\uslsh}{\strut\textbf{\znak{#1}\thezad.}\ #3\strut}}%
  \setlength{\uslvys}{\dimexpr\ht\uslbox+\dp\uslbox\relax}%
  \ifdim\uslvys<\minpole\setlength{\uslvys}{\minpole}\fi
  \par\vspace{\mezhzad}\noindent
  \begin{minipage}[t]{\polesh}\vspace{0pt}\vspace*{-\podnyat}#2\end{minipage}%
  \hfill
  \begin{minipage}[t]{\uslsh}\vspace{0pt}\usebox{\uslbox}\end{minipage}\par}

\newcommand{\shapka}{%
  \noindent{\LARGE\bfseries Один ответ, несколько, ни одного}%
  \hfill{\small Кружок, 7\,И\quad·\quad23 сентября}%
  \vspace{4pt}\hrule\vspace{10pt}
  \noindent Фамилия, имя:\ \hrulefill\rule[-3mm]{0pt}{8mm}\par
  \vspace{10pt}
  \setcounter{zad}{0}}

\begin{document}
\shapka

%=====================================================================
% 1. ЛИНИЯ А, опора. Конкретное число: тип ответа от n не зависит.
%    Цель — 100%. Ставит линию листка: ответов бывает два.
%    «Найдите все возможности» СНЯТО по замечанию владельца:
%    это подразумевается всегда и в подсказке не нуждается.
%=====================================================================
\zad{}{\poleA}{На доске написаны десять утверждений:\newline
\textit{«среди этих десяти утверждений все истинны»}, \textit{«…ровно одно ложно»},
\textit{«…ровно два ложных»}, и так далее до последнего:
\textit{«…ровно девять ложных»}.\newline
Какие из этих утверждений истинны, а какие ложны?}

%=====================================================================
% 2. ЛИНИЯ В, опора. Конкретное число: ответ «ни одного» при любом.
%    Здесь же вводятся рыцари и лжецы — один раз на весь листок.
%=====================================================================
\zad{}{\poleA}{На острове живут рыцари, которые всегда говорят правду,
и лжецы, которые всегда лгут. В шеренгу встали 30 островитян, и каждый сказал:
\textit{«слева от меня рыцарей больше, чем справа»}.\newline
Сколько рыцарей в шеренге?}

%=====================================================================
% 3. ЛИНИЯ Б, опора. Переставлена сюда из шестой позиции по оценке владельца:
%    задача простая. Пункты СНЯТЫ — «(а) без (б) бессмысленно: как понять,
%    что это максимум». Вопрос на максимум сам требует обеих половин.
%=====================================================================
\zad{}{\poleA}{В ряд стоят 100 островитян. Каждый, кроме двух крайних,
сказал: \textit{«мои соседи~--- люди разных типов»}. Крайние промолчали.\newline
Какое наибольшее число рыцарей может стоять в ряду?}

%=====================================================================
% 4. ЛИНИЯ Б, вершина. ЗДЕСЬ n НУЖЕН: ответ зависит от кратности трём.
%    Идёт сразу за рядом: та же конструкция, но круг вместо ряда
%    и параметр вместо числа.
%=====================================================================
\zad{}{\poleB{а}{б}}{За круглым столом сидят $n$ островитян. Каждый сказал:
\textit{«среди двух моих соседей ровно один рыцарь»}.\newline
\textbf{(а)}~При каких $n$ такое возможно?\newline
\textbf{(б)}~Сколько рыцарей может сидеть за столом?}

%=====================================================================
% 5. ЛИНИЯ В, вершина. Подсчёт двумя способами.
%    Нечётность — существенна, без неё получается только «не меньше».
%    Взаимность дружбы из условия НЕ УБИРАТЬ: без неё вывод переворачивается.
%=====================================================================
\zad{}{\poleA}{На острове живёт нечётное число жителей. Дружба взаимна: если
кто-то дружит с вами, то и вы дружите с ним. Каждый рыцарь заявил:
\textit{«я дружу ровно с одним лжецом»}, а каждый лжец заявил:
\textit{«я не дружу ни с одним рыцарем»}.\newline
Кого на острове больше — рыцарей или лжецов?}

%=====================================================================
% 6. ЛИНИЯ А, вершина. ЗДЕСЬ n НУЖЕН: при чётных ответ есть, при нечётных нет.
%    Пунктов с частными случаями НЕТ сознательно: частный случай
%    («не более k рыцарей», 10 островитян) класс уже решал 15.09.
%=====================================================================
\zad{}{\poleA}{На доске написаны $n$ утверждений: первое~---
\textit{«среди этих утверждений не менее одного ложного»}, второе~---
\textit{«…не менее двух ложных»}, и так далее до $n$-го~---
\textit{«…не менее $n$ ложных»}.\newline
При каких $n$ такое возможно и сколько тогда ложных утверждений?}

%=====================================================================
% 7. ЗВЁЗДОЧКА-ГРОБ. Заменила задачу про четырёх математиков: та была
%    арифметикой в логической обёртке, а нужна логика.
%    Числа в условии НЕТ сознательно: чётность здесь не дана, а доказывается.
%    Проверено перебором: при нечётном числе жителей расстановки не существует,
%    при чётном во ВСЕХ расстановках рыцарей ровно половина.
%=====================================================================
\zad{\star}{\poleA}{На острове живёт несколько жителей. Каждый из них указал
на одного другого жителя и сказал: \textit{«он лжец»}. Оказалось, что каждый
житель был назван ровно один раз.\newline
Докажите, что рыцарей и лжецов на острове поровну.}

%=====================================================================
% 8. ЗВЁЗДОЧКА-ПАРАДОКС. Необязательная, выбивается из листка.
%    Четыре пункта: (а)(б)(в) берутся по очереди, (г) почти никто.
%=====================================================================
\zad{\star}{\poleD{а}{б}{в}{г}}{Учитель сказал кружковцам: \textit{«в январе кружок
проходит 13, 17, 20, 24, 27 и 31 числа. В один из этих дней вам будет дана
контрольная работа, но накануне вы знать об этом ещё не будете»}.\newline
Докажите, что контрольная не могла быть дана \textbf{(а)}~31 января;
\textbf{(б)}~27 января; \textbf{(в)}~20 января.\newline
\textbf{(г)}~Но 20 января контрольная состоялась, и накануне о ней действительно
никто не знал. Как это совместить с пунктами \textbf{(а)}--\textbf{(в)}?}

%=====================================================================
% ЗАПАС — две задачи средней сложности.
% Обе проверены перебором, обе на идеи, которых в листке ещё нет:
% первая — разбор по чётности с ответом «так не бывает»,
% вторая — переход к дополнению множества.
% Чтобы поставить в листок, снять % и убрать две задачи из середины:
% иначе профиль долей поедет и коридор сумм разъедется.
%=====================================================================
% \zad{}{\poleA}{В комнате 25 островитян, и каждый из них сказал:
% \textit{«число рыцарей среди нас чётно»}. Сколько рыцарей в комнате?}
%
% \zad{}{\poleA}{Опросили 100 человек. Чай любят 70 из них, кофе~--- 75,
% сок~--- 80, молоко~--- 85. Докажите, что хотя бы 10 человек любят все
% четыре напитка.}

\end{document}
